
\frontpage

\mcquestion{Which instrument is most suitable to determine the volume of a small irregularly shaped stone?}
\mcqchoices
  {30 cm ruler}
  {digital timer}
  {measuring cylinder}
  {tape measure}

\mcquestion{A steel ball is dropped from the top floor of a building. Air resistance can be ignored.}
\mcqtext{Which statement describes the motion of the ball?}
\mcqchoices
  {The ball falls with constant acceleration.}
  {The ball falls with constant speed.}
  {The ball falls with decreasing speed.}
  {The ball falls with increasing acceleration.}

\mcquestion{The diagram shows a velocity-time graph for an object which is accelerating.}
\mcqimage[82mm]{\assetpath/mcq/q03_velocity_time_graph.png}
\mcqtext{What is the acceleration of the object?}
\mcqchoiceswide
  {0.40 m/s\textsuperscript{2}}
  {2.5 m/s\textsuperscript{2}}
  {3.0 m/s\textsuperscript{2}}
  {100 m/s\textsuperscript{2}}

\newpage

\mcquestion{Diagram 1 shows a piece of flexible material that contains many pockets of air. Diagram 2 shows the same piece of flexible material after it has been compressed so that its volume decreases.}
\mcqimage[105mm]{\assetpath/mcq/q04_compression_diagrams.png}
\mcqtext{What happens to the mass and to the weight of the flexible material when it is compressed?}
\begin{mcqtable}{|c|c|c|}
\hline
 & mass & weight \\ \hline
\textbf{A} & increases & increases \\ \hline
\textbf{B} & increases & no change \\ \hline
\textbf{C} & no change & increases \\ \hline
\textbf{D} & no change & no change \\ \hline
\end{mcqtable}

\mcquestion{A student determines the density of copper. She has a rectangular block of copper.}
\mcqtext{Which other apparatus does she need?}
\begin{mcqtable}{|c|c|c|c|c|}
\hline
 & balance & ruler & stop-watch & thermometer \\ \hline
\textbf{A} & \checkmark & \checkmark & \checkmark & \checkmark \\ \hline
\textbf{B} & \checkmark & \checkmark & \checkmark & \textit{x} \\ \hline
\textbf{C} & \checkmark & \checkmark & \textit{x} & \textit{x} \\ \hline
\textbf{D} & \checkmark & \textit{x} & \textit{x} & \textit{x} \\ \hline
\end{mcqtable}

\mcquestion{A 0.15 kg mass is suspended from a vertical spring. The mass is pulled downwards until the spring has extended by 12 cm from its unstretched length.}
\mcqtext{The spring constant of the spring is 30 N/m. What is the magnitude of the initial acceleration of the spring when it is released?}
\mcqchoiceswide
  {14 m/s\textsuperscript{2}}
  {24 m/s\textsuperscript{2}}
  {28 m/s\textsuperscript{2}}
  {34 m/s\textsuperscript{2}}

\newpage

\mcquestion{A cricket ball of mass 0.15 kg is moving at a speed of 30 m/s.}
\mcqtext{The diagram shows a batsman holding a bat stationary as the ball hits the bat and bounces back at the same speed in the opposite direction.}
\mcqimage[90mm]{\assetpath/mcq/q07_cricket_bat.png}
\mcqtext{The ball is in contact with the bat for $1.0 \times 10^{-3}$ s. Which row is correct?}
\begin{mcqtable}{|c|c|c|}
\hline
 & \begin{tabular}{c}change of momentum\\kg m/s\end{tabular} & force applied to ball / N \\ \hline
\textbf{A} & 4.5 & 4500 \\ \hline
\textbf{B} & 4.5 & 9000 \\ \hline
\textbf{C} & 9.0 & 4500 \\ \hline
\textbf{D} & 9.0 & 9000 \\ \hline
\end{mcqtable}

\mcquestion{A stone is thrown vertically upwards at a speed of 5.0 m/s.}
\mcqtext{All of the energy initially in the kinetic store of the stone is transferred to its gravitational potential store when the stone reaches its maximum height. Which equation gives the maximum height $h$ reached by the stone?}
\mcqchoices
  {$h=\dfrac{9.8}{2 \times 5.0^2}$}
  {$h=\dfrac{5.0^2 \times 2}{9.8}$}
  {$h=\dfrac{5.0^2}{2 \times 9.8}$}
  {$h=\dfrac{2 \times 9.8}{5.0^2}$}

\newpage

\mcquestion{Which wave will diffract the most when passing through a gap that is 0.85 m wide?}
\begin{mcqtable}{|c|c|c|}
\hline
 & \begin{tabular}{c}wave speed\\m/s\end{tabular} & wave frequency / Hz \\ \hline
\textbf{A} & 300 & 280 \\ \hline
\textbf{B} & 300 & 400 \\ \hline
\textbf{C} & 350 & 280 \\ \hline
\textbf{D} & 350 & 400 \\ \hline
\end{mcqtable}

\mcquestion{Which arrow on the graph shows the amplitude of the wave?}
\mcqimage[110mm]{\assetpath/mcq/q18_wave_amplitude.png}

\mcquestion{Which row describes a seismic S-wave?}
\begin{mcqtable}{|c|p{30mm}|p{92mm}|}
\hline
 & type of wave & direction of vibration \\ \hline
\textbf{A} & longitudinal & parallel to the direction of travel of the wavefront \\ \hline
\textbf{B} & longitudinal & perpendicular to the direction of travel of the wavefront \\ \hline
\textbf{C} & transverse & parallel to the direction of travel of the wavefront \\ \hline
\textbf{D} & transverse & perpendicular to the direction of travel of the wavefront \\ \hline
\end{mcqtable}

\newpage

\mcquestion{A ray of light, travelling in air, is incident on an air-glass boundary.}
\mcqtext{The angles between the incident ray, the refracted ray and the boundary are shown.}
\mcqimage[70mm]{\assetpath/mcq/q21_refraction.png}
\mcqtext{Which equation is used to calculate the refractive index $n$ of the glass?}
\mcqchoiceswide
  {$n=\dfrac{\sin40^\circ}{\sin26^\circ}$}
  {$n=\dfrac{\sin40^\circ}{\sin64^\circ}$}
  {$n=\dfrac{\sin50^\circ}{\sin26^\circ}$}
  {$n=\dfrac{\sin50^\circ}{\sin64^\circ}$}

\mcquestion{A small object O is placed near a converging lens, as shown. The lens forms an image I.}
\mcqimage[95mm]{\assetpath/mcq/q22_lens_image.png}
\mcqtext{Which statement is correct?}
\mcqchoices
  {The image I is diminished.}
  {The image I is inverted.}
  {The image I is real.}
  {The object O is closer to the lens than its principal focus.}

\newpage

\mcquestion{A metallic wire X has a resistance.}
\mcqtext{What is the difference in resistance for a wire of the same material as X that is shorter than X but the same diameter, and for a wire that has the same length as X but a smaller diameter?}
\begin{mcqtable}{|c|p{42mm}|p{48mm}|}
\hline
 & resistance of shorter wire & resistance of wire with smaller diameter \\ \hline
\textbf{A} & greater & greater \\ \hline
\textbf{B} & greater & smaller \\ \hline
\textbf{C} & smaller & greater \\ \hline
\textbf{D} & smaller & smaller \\ \hline
\end{mcqtable}

\mcquestion{The diagram shows the circuit for a potential divider.}
\mcqimage[56mm]{\assetpath/mcq/q29_potential_divider.png}
\mcqtext{Which equation is used to calculate $V_1$?}
\mcqchoices
  {$V_1=\dfrac{R_1}{V_2 \times (R_1+R_2)}$}
  {$V_1=\dfrac{V_2 \times R_1}{R_2}$}
  {$V_1=\dfrac{V_2 \times R_1}{R_1+R_2}$}
  {$V_1=\dfrac{R_1}{V_2 \times R_2}$}

\newpage

\mcquestion{When there is an electric current in a long straight wire, a magnetic field is created around the wire.}
\mcqtext{Which diagram shows the correct pattern and direction of magnetic field lines around a long straight wire carrying current into the page?}
\mcqimage[118mm]{\assetpath/mcq/q32_magnetic_field_choices.png}

\mcquestion{When a thin gold foil is bombarded with alpha particles, some alpha particles are deflected through large angles.}
\mcqtext{Which statement explains this deflection?}
\mcqchoices
  {Most of the atom consists of empty space.}
  {All of the positive charge and most of the mass of the gold atom are concentrated in a small volume.}
  {Positive charge in the gold atom is spread evenly throughout the atom.}
  {All of the negative charge is concentrated at its centre.}

